\chapter*{\uppercase{DECLARATION ON THE USE OF AI TOOLS}}
\addcontentsline{toc}{chapter}{\bfseries \uppercase{DECLARATION ON THE USE OF AI TOOLS}}
The memory engine uses AI models as parts of the system, to create its synthetic training data
and in its evaluation. We designed the system, checked the training labels, ran and checked
every experiment, and are responsible for all the content of this report.

The following AI tools were used:

\textbf{Tool name(s):} Qwen3-4B-Instruct-2507 (and Qwen3-1.7B in early trials);
\code{qwen3-embedding:0.6b}; GPT-5.6-sol; Claude Sonnet 5.5; Claude Haiku 5.5; TypeSafe Jev and
GPT-4o-mini.

\textbf{Purpose of use:} base model for our fine-tuned extractor \code{mem-extractor}, which
runs M2 extraction, consolidation and reflect answers (Qwen3-4B-Instruct-2507); the embedding
model \code{qwen3-embedding:0.6b}, which turns text into vectors for search; writing and labelling the synthetic training
conversations (GPT-5.6-sol) and labelling the coverage set (Claude Sonnet 5.5); the Pi agent's
model in the live demo (Claude Sonnet 5.5) and its reader model in the benchmark (Claude Haiku 5.5); benchmark judges (TypeSafe Jev,
GPT-4o-mini).

\textbf{Extent of use:} all training data is synthetic and was generated by the models named
above. The extractor and embedding models run locally through Ollama. The reader and judge
models were used only to produce and grade benchmark answers. All measured results in
Chapter~\ref{ch:results} come from our own runs.
